\scanpage{50}
\mspage{48}

alors on peut la réaliser par une flèche \uline{dans Cat} qui rende le diagramme comm. pour de bon (\unsure{dans} Cat). On est ramené au $=$ problème dans une \deleted{\unsure{situa.}}\unsure{situation}

\diagram{\begin{tikzcd}[ampersand replacement=\&]
X \arrow[r,hook,"i"'] \& Y \arrow[l,bend right=42,dashed,"p"']
\end{tikzcd}},

où $X$ est une catégorie de Thomason, $i$ une cofibration de Thomason, $p$ une rétraction de $Y$ dans $X$ \uline{dans} \uline{Cat hot}, on veut la réaliser par une vraie rétraction.

\margin{? $\big|$}

En admettant que ce soit le cas, donc que \deleted{\unsure{tout}} hom \deleted{\unsure{des}} ind objets $\underline{\underline K},\underline{\underline {K'}}\deleted{\unsure{{}_{\infty}}}$ \uline{dans Cat hot} se représente \unsure{par} un hom \deleted{de} entre les ind-objets \uline{dans Cat}, donc aussi par un foncteur $\deleted{\unsure{C}}K_\infty\to K'_\infty$, reste la question si cet homomorphisme \inserted{correspondant dans Cat hot} est unique. \deleted{dans le}

Donc si on a deux \deleted{\unsure{hom}} foncteurs $f,g$ qui induisent les $=$ hom sur les foncteurs ind-représentés \unsure{p.} sur Cat hot, il faut \unsure{q|que} \unsure{soient} homotopes. \inserted{Donc} \deleted{\unsure{on a donc à cons.}}

% slanted
\margin{$K,K'$ catégories de Thomason, les $K_n,K'_n$ aussi}

\uline{Question} Deux foncteurs $f,g\colon K_\infty\unsure{\relbar\joinrel\relbar|\longrightarrow} K'_\infty$ \deleted{\unsure{qui}} tels, que pour $\forall n\in\mathbb N$, \unsure{$f_n,g_n$|$f,g$} qui \unsure{s'en} déduisent par restriction à $K_n$ sont homotopes \unsure{?|,} (et $=$ homotopes en tant que foncteurs de $K_n$ dans un $K'_N$ convenable) \deleted{\unsure{que}} sont-ils homotopes ?

% archive sheet number: 49
